\scanpage{152}
\mspage{5}
% archive foliation: 151

1°) Cas où $X$ équi-singulier autour de $Y$.
Je présume qu'on doit pouvoir prendre
$\varphi$ fibrant (``rigidification fibrante''),
et qu'on a une équivalence de catégories
$\infty$-isotopiques

\diagram{\begin{tikzcd}[column sep=small,row sep=large]
& \text{Syst. }(X',B,\deleted{\varphi}\inserted{Y},\varphi)
  \arrow[d,"\approx"]
& \begin{gathered}(X',B)\text{ paire bordante}\\
\varphi:B\longrightarrow Y\text{ morphisme}\\\text{fibrant}\end{gathered}\\
\begin{gathered}(X',B,Y,\varphi)\\
\longmapsto((X'\amalg_B Y,\varphi),Y)\end{gathered}
& (X,Y):\text{ paires équi-singulières.} &
\end{tikzcd}}

\uline{NB} les $X\setminus Y,Y$
\deleted{\unsure{non singuliers, alors}}
\inserted{réguliers}
doivent correspondre au cas où
$X'$ variété à bord $B$, et $Y$ régulier.

\margin{Je présume que $X\setminus Y$ régulier
implique que $Y$ régulier
(au moins si $Y\subset\overline{X\setminus Y}$)}

2°) Cas général (pas néc. équi-sing.).
Il faudrait arriver à imposer des
conditions suppl. sur $\varphi:\dot V\longrightarrow Y$,
pour assurer l'unicité \unsure{esselle}.
Peut être qque chose du genre
suivant : Pour $\unsure{\forall}\ y\in Y$, si $c(y)=\operatorname{codim}_y(Y,X)$
alors $\dim\varphi^{-1}(y)=c(y)-1$
(\unsure{cas de figure} plus précis, local sur $\dot V$ \unsure{ald} de l'être sur $Y$
